% !TeX root = ../papers/nc-p3.tex
\section{平面商の全Jordan型と平面を忘れた同型}
\label{sec:plane-quotient-normal-forms}
正則正規一次直線 $kt$ を中央化し、三生成商が平面グラフとなる枝を扱う。
$H\in\mathrm{GL}_3$ を射影共役まで取り、巡回添字 $(i,j,k)$ に対し
$r_i=x_j H(x_k)-x_k H(x_j)$ を商の関係とする。
四生成関係は $t$ の中心性と $r_i=t\sum_j M_{ji}x_j+t^2 b_i$ である。
以下の表で与える全係数が標準Koszul左右AS正則、最小周期1である。
\begin{theorem}\label{thm:plane-quotient-all-marked}
一つの平面を指定した添字保持同型類は、以下の各型で尽きる。
\begin{enumerate}
\item $H$ が非スカラーで対角化でき、固有値を $h_1,h_2,h_3$ とする。
$h_i^3=\det H$ を満たす添字だけに $M_{ii}=a_i,b_i$ を許し、他の係数を零とする。
該当添字数は0、1、3である。一つなら二係数の非零台だけを記録する。
三つなら比が原始三乗根の固有値になり、関係左辺を同じ巡回交換係数で規格化する。
この場合は六係数の台とトーラス
$a_i\mapsto d_i a_i/(d_jd_k),\ b_i\mapsto b_i/(d_jd_k)$ の整数核単項式を記録し、
循環置換も除く。
\item $H=J_2(1)\oplus[\lambda]$, $\lambda\ne1$ のとき、
$\lambda\ne-1$ では補正は全て零である。$\lambda=-1$ では $M_{33},b_3$ の二係数の台だけを記録する。
\item $H=J_3(1)$ のとき
\[
 M=\begin{pmatrix}0&0&0\\2a/3&-a/3&0\\a&0&0\end{pmatrix},
 \qquad b=(b_0,0,a^2/9)^t .
\]
$(a,b_0)$ の代表は $(1,0),(0,1),(0,0)$ の三つである。
\item $H=J_2(1)\oplus[1]$ では
\[
 M=\begin{pmatrix}0&-a/2&0\\a&0&-2c\\b&c&d\end{pmatrix},\qquad
 (b_0,b_1,b_2)^t\quad\text{を定数欄とする。}
\]
標準代表を次に示す。最初の二行の定数は全て零であり、その後の行では $b_1=0$。
\[
\begin{array}{c|c|l}
 \text{条件}&(a,b,c,d)&(b_0,b_2)\\ \hline
 c\ne0 &(0,-1/2,1,0)&(0,0)\\
 c=0,\ a\ne0 &(1,0,0,0)&(0,0)\\
 a=c=0,\ d\ne0 &(0,0,0,1)&(0,0)\text{ または }\mathbb P(3,2)\\
 a=c=d=0,\ b\ne0 &(0,1,0,0)&(0,v)\ (v\in k),\quad(1,1)\\
 M=0 &(0,0,0,0)&(0,0),(1,0),(0,1)
\end{array}
\]
\item $H$ がスカラーなら、三次元Lie環 $\mathfrak g$ とそのスカラー二コサイクル類の表を使う。
アーベル、Heisenberg、$\mathfrak{aff}_1\oplus k$ ではそれぞれ零類と一つの非零類である。
二次元アーベル導来環に残る一文字が可逆行列 $P$ で作用する型は $[P]$ の共役類で分類し、
跡零のときに限り零類と非零類に分かれる。$\mathfrak{sl}_2$ は零類だけである。
\end{enumerate}
\end{theorem}
\begin{proof}[PBW条件の座標化]
$R_i$ を $r_i$ の二文字係数行列、$D=\det H$ と置く。
三次交叉は交代体積に $1\otimes H\otimes H^2$ を作用させたもので、
右端の文字は $H^2x_i$、左端の文字は $DH^{-1}x_i$ である。
従って定理\ref{thm:regular-normal-linear}の条件は
\[
 F=(M(H^2)^t-DH^{-1}M^t)H^{-t},\quad F+F^t=0,\quad
 C=(F_{23},F_{31},F_{12})^t,\qquad
 MC+(H^2-DH^{-1})b=0,\quad b^t C=0
\]
である。持上げ変更 $e$ は $M$ の第 $i$ 列から $(R_i+R_i^t)e$ を引き、
定数列に $M^te-(e^tR_i e)_i$ を加える。
$H$ と可換する基底変換 $G$ は持上げ変更の前に
$(M,b)\mapsto(GMG^t/\det G,Gb/\det G)$ と作用する。
これは $r_i$ の変更行列が $\det(G)G^{-t}$ であるためである。
以上は各Jordan行列についても分母の消滅を伴わない有限計算である。

非スカラー対角の場合、$F+F^t=0$ の非対角式は
$(h_j-h_k)M_{ij}+(h_i-h_k)M_{ji}=0$ である。一対ごとに二係数の差の一方は非零で、
残る一方向は対応する持上げ変更で消せる。対角式は $(h_i^3-D)M_{ii}=0$ となる。
このとき $C=0$ なので定数にも同じ共鳴条件が掛かる。
固有値が重複し共鳴が一文字だけのとき、その二係数は外の文字の倍率と
二文字ブロックの行列式で独立に変えられる。相異なる固有値の場合にも同じく二倍率が独立である。
全三文字が共鳴する非スカラーの場合は固有値が三乗根の一軌道になり、固定した対角表示を
射影的に保つ有限置換は三巡回だけである。

$J_2(1)\oplus[\lambda]$ では、中央化群は第一ブロックの上三角Toeplitz行列と
第三文字の倍率である。先の線型式を消去すると $\lambda\ne\pm1$ では全 $M$ を消せ、
$H^2-DH^{-1}$ も可逆である。$\lambda=-1$ ではただ一つの $M_{33}$ と $b_3$ が残り、
ブロック倍率の二乗と第三倍率が台上で推移的に働く。
$J_3(1)$ では持上げ変更後の式は表の通りで、$G=pI+q(H-I)+r(H-I)^2$ の作用で
$a\mapsto a/p$ となる。$a\ne0$ なら $b_0$ に対する $r$ の一次係数は
$-8a^2/(9p^3)$ で非零、$a=0$ なら $b_0\mapsto b_0/p^2$ である。
\end{proof}

\begin{proof}[二重Jordanの残余作用とLie型]
$H=J_2(1)\oplus[1]$ で $F$ の線型式と持上げ変更を使うと定理の $M$ となる。
定数についての式は
\[
 3ad+4bc+2c^2=0,\qquad b_1=a(b-c)/6,\qquad
 -3a b_2+2(b-c)b_1-6c b_0=0 .
\]
未使用の第二持上げ変更は定数に $(a,0,-2c)^t$ の任意の倍を加える。
$G$ は行を $(p,q,r),(0,p,0),(0,z,t)$ と書ける（$pt\ne0$）。
$M$ を再び表の形に戻すと
\begin{align*}
 c'&=c/p,& a'&=(ap-2cr)/(pt),& d'&=(dt-cz)/p^2,\\
 b'&=(3apz+4bpt-6crz+6drt)/(4p^2t).
\end{align*}
$c\ne0$ なら $c'=1$, $a'=d'=0$ とでき、方程式が $b'=-1/2$ を与える。
$a\ne0=c$ なら $d=0$ で $a'=1,b'=0$ とできる。いずれも残った定数は
未使用の変更で全て零になる。$a=c=0$ なら $d\ne0$ で $d'=1,b'=0$ とでき、
この表示の残余作用は $(b_0,b_2)\mapsto (b_0/p^3,b_2/p^2)$ である。
$d=0,b\ne0$ では $b=1$ に直した後は $(b_0,b_2)\mapsto
((b_0+r(b_2-1))/t,b_2)$。$M=0$ の場合は
$(b_0,b_2)\mapsto(b_0/(pt)+r b_2/(p^2 t),b_2/p^2)$ となる。
これで全ての零値を含め表の代表が得られる。
\path{computations/verify_plane_quotient_jordan.py} は線型消去、残余二次式、倍率を記号的に照合する。

スカラーの場合、$\alpha:\bigwedge^2 V\to V$ はJacobi恒等式を満たす括弧、
定数は $d\beta=0$ を満たし、持上げ変更はコバウンダリの加算である。
導来環の次元0では交代形式の階数が0か2。次元1では
$[x,y]=z$（中心導来）または $[h,x]=x$（非中心導来）に基底を取れる。
前者の $H^2$ は $x^*\wedge z^*,y^*\wedge z^*$、後者は
$h^*\wedge z^*$ で張られ、自己同型は非零類で推移的に働く。
次元2では可解Lie環の導来環が冪零であることから二次元アーベルであり、
残る文字の作用 $P$ は全射である。コサイクル式は
$\operatorname{tr}P\cdot\beta(x,y)=0$、残る二成分はコバウンダリである。
跡零なら $P$ の中央化群の行列式で非零類を同定できる。
導来次元3では可解でなく、Levi分解と次元から $\mathfrak{sl}_2$ である。
括弧が $\bigwedge^2V\to V$ の同型なので全ての定数はコバウンダリとなる。
\end{proof}

\begin{theorem}\label{thm:plane-quotient-unmarked}
上の全表について平面を指定しない同型類は、可換多項式窓を除いて、
以下で得る高々二つの指定付き標準形の順序なしリストで決まる。
最初の平面 $t=0$ に加え、四つの可換変数での三つの二次式
\[
 f_i(x,t)=r_i(x,x)-t\sum_j M_{ji}x_j-t^2b_i
\]
の共通一次因子 $v\notin kt$ があれば、その平面を第二の印に取る。
各印について定理\ref{thm:finite-normal-line-recognition}の $g_i$ で中央化し、
表の該当する代表を記録する。可換多項式窓は一つの同型類である。
これらは平面商の全特殊値に通用し、全窓の各辺射影ファイバーは空または一点である。
\end{theorem}
\begin{proof}
$t=0$ の点で中心関係は次の $t'=0$ を強制し、次点は三生成商の射影自己同型で一意である。
$t\ne0$ では中心関係により次点は同じ点になり、残る条件は $f_i=0$ である。
右からも同じなので全射影に正次元ファイバーはない。
外に開部分を持つ別の平面は従って共通一次因子に等しく、反対に各その因子は
対角グラフを与える。全 $f_i=0$ なら三生成関係は交代関係で両補正も零である。
それ以外では命題\ref{prop:surface-degree-bound}から平面は合計高々二つである。
全位置で同じ表示を用いたとき最初の平面はそのグラフで移り、第二平面は対角で移る。
定理\ref{thm:plane-surface-normal}から両者とも正規一次直線列を与える。
同型は平面を送るのでリストは不変であり、一つの印を合わせて標準形が一致すれば
定理\ref{thm:regular-normal-linear}が裸の窓の同型も与える。
\end{proof}
